
% File hicss.tex
%
% Contact: Holm Smidt, hicss@hawaii.edu
%%
%%
%% Based on the style files for ACL 2015 by 
%% car@ir.hit.edu.cn, gdzhou@suda.edu.cn


\documentclass[10pt]{article}
\usepackage[letterpaper]{geometry}
\usepackage{hicss}
\usepackage{times}
\usepackage[none]{hyphenat}
\usepackage{url}
\usepackage{latexsym}
\usepackage{indentfirst}
\usepackage{graphicx}
\graphicspath{{images/}}
\usepackage{amsmath}
\usepackage{booktabs}
\usepackage{xcolor}
\usepackage{enumitem}
\setlist{noitemsep}
\usepackage{csquotes}
\usepackage[
    style=apa,
  ]{biblatex}
\addbibresource{references.bib}

\newcommand{\sansserifformat}[1]{\fontfamily{cmss}{ #1}}%

\setlength\titlebox{7cm}

% title for manuscript without author info 
\title{Discovering Neurobiological Subtypes in Autism via Dynamic Spatiotemporal Graph Attention Networks}

% Comment this for initial manuscript 
% Uncomment this for final manuscript
% \author{Aadi Sharma \\
%  SRM Institute of Science and Technology \\
%  {\underline{as2382@srmist.edu.in}} \\ \\ WRITE ADI EMAIL HERE
%  Ashmita Das \\
%  SRM Institute of Science and Technology \\
%  {\underline{ad8902@srmist.edu.in}} \\ \And           WRITE ASHMITA HERE
%  Suryansh Maurya \\
%  SRM Institute of Science and Technology \\
%  {\underline{sm0464@srmist.edu.in}} \\ \\  SUYANSH HERE
%  Abinaya Gopalakrishnan \\
%  SRM Institute of Science and Technology \\
%  {\underline{abinayag2@srmist.edu.in}} \\   ABHINYA HERE yess sahi hai mere alawa baki ke
%  Hrishitva Patel \\
%  University of Texas at San Antonio \\
%  {\underline{hrishitva.patel@utsa.edu}} \\ }

\date{}

\begin{document}
\maketitle
\begin{abstract}
Autism Spectrum Disorder (ASD) is a very complex  neurodevelopmental condition which is characterized by significant heterogeneity. Neuroimaging research has identify consistent biomarkers, but efforts have been hampered by the dynamic nature of brain function and high inter-subject variability. We propose the Dynamic Spatiotemporal Graph Attention Transformer (DST-GAT), a novel geometric deep learning framework designed to explicitly model the time-varying nature of brain networks. For each subject's rs-fMRI scan from the ABIDE I dataset, we generate a sequence of dynamic functional connectivity graphs using sliding window approach on time-series extracted from the AAL116 atlas. The proposed DST-GAT framework integrates a Graph Attention Network (GAT) to capture spatial dependencies within the brain's connectome topology, alongside a Temporal Transformer to model temporal dynamics across brain states. This combination results in a comprehensive attention-based architecture capable of identifying critical spatiotemporal patterns within neural data. To ensure robustness, we used 5-fold cross-validation strategy complemented by advanced regularization techniques, including Stochastic Weight Averaging (SWA), for consistent and unbiased model evaluation. The DST-GAT achieved a mean Area Under the Curve (AUC) of 0.74 ± 0.04 across validation folds, outperforming all evaluated baseline models. Beyond classification, the model was further utilized as a feature representation extractor. Through unsupervised clustering (using UMAP for di and HDBSCAN), we identified three distinct neurobiological subgroups within the ASD population. Subsequent post-hoc analyses revealed significant connectivity variations, primarily within the default mode, salience, and fronto-parietal networks. By incorporating the dynamic characteristics of brain connectivity, The proposed framework demonstrates how spatiotemporal graph modeling can support ASD classification and data-driven subgroup analysis. The emergence of data-driven neurobiological subtypes highlights its potential to address ASD's clinical heterogeneity and advance personalized treatment approaches. 
\end{abstract}

\subsubsection*{Keywords:}

Autism Spectrum Disorder, rs-fMRI, Graph Neural Networks, Spatiotemporal Modeling, Deep Learning, DST-GAT

% --- SECTION 1: INTRODUCTION ---
\section{Introduction}
Autism Spectrum Disorder (ASD) is marked by substantial clinical and neurobiological heterogeneity. Individuals sharing the same diagnosis often exhibit distinct behavioral and neural profiles, suggesting multiple underlying mechanisms.
One of the central challenge in understanding ASD, both clinically and scientifically, is its remarkable heterogeneity \parencite{apa2013dsm5, lombardo2019autism}. Because ASD exists on a spectrum, individuals who share the same diagnostic label may exhibit vastly different symptom patterns and severity levels, implying that distinct neurobiological mechanisms could underlie the disorder.

Resting-state fMRI provides a non-invasive view of large-scale brain
network organization and has been widely used in ASD research \parencite{biswal1995functional}.
However, most prior studies rely on static connectivity representations,
which fail to capture the dynamic nature of neural interactions \parencite{hutchison2013dynamic}. This
makes the use of spatiotemporal models which can explicitly learn
time-varying patterns in brain connectivity.


Our contributions are threefold:
\begin{enumerate}
    \item \textbf{A Novel Architecture:} We present the DST-GAT, a fully attentional model that synergistically combines a Graph Attention Network (GAT) for learning spatial graph features with a Temporal Transformer for learning long-range dependencies between graph states. This allows the model to identify which brain connections are most important at each moment (spatial attention) and which moments in the scan are most informative for the final diagnosis (temporal attention).
    \item \textbf{State-of-the-Art Performance:} The large-scale ABIDE I dataset that the DST-GAT significantly outperforms traditional static models and simpler recurrent spatiotemporal models, achieving a new benchmark in ASD classification from rs-fMRI data.
    \item \textbf{Discovery of Neurobiological Subtypes:} We pioneer a ``classification-for-discovery'' approach. By using the trained DST-GAT as a powerful feature extractor, we project subjects into a learned embedding space. Unsupervised clustering within this space reveals three distinct and robust neurobiological subtypes within the ASD cohort, providing a data-driven path to stratifying the autism spectrum.
\end{enumerate}

% --- SECTION 2: RELATED WORK ---
\section{Related Work}
The uses of machine learning to rs-fMRI data for ASD classification has a rich history, evolving from classical methods to sophisticated deep learning architectures. This section reviews the major trends and identifies the gaps that motivate our work.

\subsection{Classical Machine Learning and Static Connectivity}
Early research on ASD classification primarily relied on traditional machine learning approaches applied to features extracted from static functional connectivity (FC) matrices. Among these, Support Vector Machines (SVMs) were the most widely adopted, often paired with feature selection strategies to mitigate the challenges posed by the high dimensionality of FC data \parencite{nielsen2013multi, heinsfeld2018identification}. Other algorithms, such as Random Forests and Logistic Regression, were also explored in similar contexts. Although these studies played an important role in establishing the feasibility of automated ASD diagnosis, their performance was constrained by limited dataset sizes and poor generalization across different imaging sites, a persistent issue particularly evident in the ABIDE dataset \parencite{di2014autism}. Moreover, their exclusive dependence on static FC features restricted their capacity to represent the temporal variability inherent in brain function, thereby overlooking crucial dynamic neural interactions that may underlie ASD \parencite{yahata2016prediction, khosla2019detecting}.

\subsection{Graph Neural Networks for Brain Connectomes}
Recognizing that the brain is a network, researchers began applying Graph Neural Networks (GNNs), particularly Graph Convolutional Networks (GCNs), to model the brain connectome \parencite{kips2017semi}. GCNs are well-suited to learn node-level and graph-level representations from static brain graphs, where the regions of brain are nodes and FC values are edge weights \parencite{parisot2018disease}. This marked a significant conceptual advance over CNN-based methods. However, most GCN applications in neuroscience have still been limited to static graphs, thus inheriting the limitations of the underlying data representation. They model the spatial topology of the connectome but not its temporal evolution \parencite{allen2014tracking, gadgil2020spatio}.

\subsection{The Emergence of Transformers}
The Transformer architecture, with its self-attention mechanism, has revolutionized sequence modeling by effectively capturing long-range dependencies without the sequential processing constraints of RNNs \parencite{vaswani2017attention}. While its initial success was in natural language processing, it is increasingly being applied to other domains, including computer vision \parencite{dosovitskiy2020image} and, more recently, time-series analysis. Our work is situated at the confluence of these advanced architectural trends. We replace the isotropic GCN with a more expressive Graph Attention Network (GAT) \parencite{velivckovic2018graph}, which apply attention to learn the relative importance of different brain connections. We then replace the LSTM with a Temporal Transformer, which uses self-attention to identify the most salient moments in the dynamic evolution of the brain connectome. This creates a powerful, end-to-end, fully attentional framework, the DST-GAT, that is purpose-built to address the challenges of learning from dynamic spatiotemporal brain data. 

% --- SECTION 3: MATERIALS AND METHODS ---
\section{Materials and Methods}

\subsection{Dataset and Subject Selection}
We used resting-state fMRI data from the ABIDE I consortium \parencite{di2014autism}. Preprocessed scans from the C-PAC pipeline were filtered using standard quality-control criteria, resulting in a final cohort of 871 subjects (418 ASD, 453 TD).

\subsection{Dynamic Graph Construction Pipeline}
The core of our methodology is the transformation of each subject's 4D rs-fMRI scan into a sequence of dynamic graphs. This pipeline is illustrated in Fig.~\ref{fig:pipeline} and detailed below.

\begin{figure}[!t]
  \centering
  \includegraphics[width=\columnwidth]{Image1.jpg}
  \caption{The data preprocessing pipeline. For each subject, mean BOLD time-series are extracted from 116 regions of interest defined by the AAL116 atlas. A sliding window (size=50, step=5) is applied to these time-series. Within each window, a Pearson correlation matrix is computed, forming the adjacency matrix of a graph for that time interval. The final output for one subject is a sequence of graphs representing the brain's dynamic functional connectivity.}
  \label{fig:pipeline}
\end{figure}

Each subject's rs-fMRI time series was transformed into a sequence of dynamic functional connectivity graphs using an overlapping sliding window strategy \parencite{hutchison2013dynamic, allen2014tracking}. Brain regions defined by the AAL116 atlas \parencite{tzourio2002automated} were treated as nodes, while time-varying Pearson correlations served as edge weights. This process yields temporally ordered graphs with fixed topology and dynamic connectivity.


\subsection{The DST-GAT Architecture}
The DST-GAT model is designed to process these sequences of graphs. It consists main components: a spatial GAT encoder, a temporal Transformer encoder, and final classification head. The overall architecture is depicted in Fig.~\ref{fig:architecture}.

\begin{figure}[!t]
  \centering
  \includegraphics[width=\columnwidth]{Image2.jpg}
  \caption{The Dynamic Spatiotemporal Graph Attention Transformer (DST-GAT) architecture. The model first uses a shared-weight GAT to encode each graph in the temporal sequence into a fixed-size embedding vector. These vectors are then augmented with positional encoding and processed by a Transformer encoder to model temporal dependencies. The final representation is passed to a simple MLP for classification.}
  \label{fig:architecture}
\end{figure}

\subsubsection{Component A: Graph Attention Network (GAT) Spatial Encoder}

\begin{equation}
e_{ij} = \text{LeakyReLU}(\mathbf{a}^T[\mathbf{W}\mathbf{h}_i \| \mathbf{W}\mathbf{h}_j])
\end{equation}
 Multi-head attention is used to stabilize learning, where multiple independent attention mechanisms are computed and their outputs are concatenated \parencite{velivckovic2018graph}. Our GAT encoder consists of 3 \texttt{GATv2Conv} layers (an improved version of the original GAT that addresses the issue of static attention and has been shown to be more expressive) with 4 attention heads each. The input node features are an identity matrix $\mathbf{X} \in \mathbb{R}^{N \times N}$, allowing the model to learn unique embeddings for each ROI. After the final GAT layer, we apply global mean pooling to convert the node features into a single embedding vector $\mathbf{z}_k \in \mathbb{R}^{D}$, where $D$ is the hidden dimension. This process is applied to every graph in the sequence, producing a sequence of embeddings $\mathbf{Z} = (\mathbf{z}_1, \mathbf{z}_2, \ldots, \mathbf{z}_K)$.

\subsubsection{Component B: Temporal Transformer Encoder and Classifier}
The sequence of graph-level embeddings $\mathbf{Z} = (\mathbf{z}_1, \mathbf{z}_2, \ldots, \mathbf{z}_K)$ produced by the GAT encoder is treated as an input sequence for a standard Transformer encoder \parencite{vaswani2017attention}. Since the Transformer architecture is permutation-invariant, we add learnable sinusoidal positional encodings to each embedding so that the model can distinguish temporal order. A special learnable \texttt{[CLS]} token is prepended to the sequence. The augmented sequence is then processed by a stack of 4 Transformer encoder layers, each consisting of multi-head self-attention (8 heads) followed by a position-wise feed-forward network with GELU activation. Through self-attention, each time step can attend to every other time step, allowing the model to capture long-range temporal dependencies across brain states and to selectively emphasize the most diagnostically informative moments in the scan.

The output corresponding to the \texttt{[CLS]} token is used as the final, fixed-size representation of the entire 4D fMRI scan. This aggregated vector captures both the spatial topology learned by the GAT and the temporal dynamics modeled by the Transformer. It is passed through a two-layer Multi-Layer Perceptron (MLP) with a hidden dimension of 128, GELU activation, and dropout (rate 0.3) between the layers. The final layer projects to a single logit, and the binary prediction (ASD vs.\ TD) is obtained by applying a sigmoid function during inference.

\subsection{Training and Evaluation Protocol}
\subsubsection{Cross-Validation Strategy}
To obtain a robust and unbiased estimate of model performance, we employed a 5-fold stratified cross-validation strategy. We split dataset into 5 folds, making sure that the proportion of ASD and TD subjects was maintained in every fold. The model was trained 5 times, with each fold being used as the validation set once. The final reported metrics are mean and standard deviation across the 5 folds.

\subsubsection{Implementation Details}
The model is implemented in PyTorch and PyTorch Geometric. We used the following hyperparameters and settings:
\begin{itemize}
    \item \textbf{Loss Function:} \texttt{BCEWithLogitsLoss}, combines Sigmoid layer and Binary Cross-Entropy loss for numerical stability.
    \item \textbf{Optimizer:} AdamW optimizer with learning rate $1 \times 10^{-4}$ and weight decay (L2 regularization) of $1 \times 10^{-5}$.
    \item \textbf{Scheduler:} We used a Cosine Annealing learning rate scheduler to smoothly decay the learning rate over the course of training.
    \item \textbf{Training Epochs:} The model was trained for maximum of 150 epochs.
    \item \textbf{Batch Size:} Batch size of 16 was used.
    \item \textbf{Regularization:} In addition to weight decay, dropout with a rate of 0.3 was applied in the GAT layers, the Transformer layers, and the final MLP classifier.
    \item \textbf{Early Stopping:} We monitored the validation AUC and stopped training if it did not improve for a patience of 20 epochs. The model weights with best validation AUC were saved.
\end{itemize}

\subsection{Unsupervised Subtype Discovery}
After the cross-validation was complete, we trained a final DST-GAT model on the entire dataset for 100 epochs using Stochastic Weight Averaging \parencite{izmailov2018averaging}. This model was then used as a feature extractor.
\begin{enumerate}
    \item \textbf{Embedding Extraction:} We passed all subjects (both ASD and TD) through the trained model and extracted the final high-dimensional embedding vector for each subject (the output of the Transformer corresponding to the \texttt{[CLS]} token, before the classifier).
    \item \textbf{Clustering:} We isolated the high-dimensional embeddings of the 418 ASD subjects. We applied the Hierarchical Density Based Spatial Clustering of Applications with Noise (HDBSCAN) algorithm \parencite{campello2013density} \textit{directly to these original, high-dimensional embeddings}. HDBSCAN is advantageous as it is not required for specifying the number of clusters beforehand and can identify noise points that do not belong to any cluster.
    \item \textbf{Visualization:} After identifying the cluster labels from HDBSCAN, we \textit{then} applied the Uniform Manifold Approximation and Projection (UMAP) algorithm \parencite{mcinnes2018umap} to the high-dimensional embeddings to project them into a 2D space \textit{for visualization purposes only}. This allows us to visualize the cluster separation found in the high-dimensional space.
\end{enumerate}
The resulting clusters represent potential data-driven neurobiological subtypes of ASD. We performed post-hoc analysis to characterize the connectivity patterns associated with each subtype.

% --- SECTION 4: RESULTS ---
\section{Results}
\subsection{Classification Performance}
We measured the performance of our proposed DST-GAT model against four relevant baselines to provide a robust comparison:
\begin{itemize}
    \item \textbf{SVM (Static FC):}  Support Vector Machine with RBF kernel applied to the flattened, time-averaged static functional connectivity matrices. This represents a standard classical machine learning approach.
    \item \textbf{3D CNN (ResNet-18):} A standard 3D ResNet-18 applied to time-averaged fMRI volumes, representing the static, non-graph-based deep learning approach. This baseline was carefully hyperparameter-tuned.
    \item \textbf{Static GCN (Parisot et al.):} A Graph Convolutional Network applied to the static FC matrices, replicating the methodology of \textcite{parisot2018disease}. This represents the static graph-based approach.
    \item \textbf{GCN-LSTM:} A spatiotemporal model that uses a GCN to encode each graph and an LSTM to process the sequence, representing a simpler dynamic graph modeling approach.
\end{itemize}

The classification results from the 5-fold cross-validation is summarized in Table~\ref{tab:results}. The DST-GAT model significantly outperformed all baselines across all evaluated metrics.

\begin{table}[htbp]
\caption{5-Fold Cross-Validation Classification Performance.}
\label{tab:results}
\begin{center}
\resizebox{\columnwidth}{!}{%
\begin{tabular}{@{}lcccc@{}}
\toprule
\textbf{Model} & \textbf{AUC} & \textbf{Accuracy} & \textbf{F1-Score} & \textbf{Precision} \\ \midrule
SVM (Static FC) & 0.65 $\pm$ 0.05 & 0.62 $\pm$ 0.04 & 0.61 $\pm$ 0.05 & 0.63 $\pm$ 0.04 \\
3D CNN & 0.62 $\pm$ 0.04 & 0.60 $\pm$ 0.03 & 0.58 $\pm$ 0.05 & 0.61 $\pm$ 0.04 \\
Static GCN & 0.67 $\pm$ 0.06 & 0.64 $\pm$ 0.05 & 0.64 $\pm$ 0.05 & 0.65 $\pm$ 0.05 \\
GCN-LSTM & 0.68 $\pm$ 0.05 & 0.65 $\pm$ 0.04 & 0.66 $\pm$ 0.05 & 0.65 $\pm$ 0.06 \\
\textbf{DST-GAT (Ours)} & \textbf{0.74 $\pm$ 0.04} & \textbf{0.70 $\pm$ 0.04} & \textbf{0.71 $\pm$ 0.05} & \textbf{0.72 $\pm$ 0.05} \\ \bottomrule
\end{tabular}%
}
\end{center}
\small{Metrics are reported as mean $\pm$ standard deviation across the 5 validation folds.}
\end{table}
Statistical significance between DST-GAT and baseline models was assessed using a paired two-tailed t-test on fold-wise AUC scores. The improvement is over the strongest baseline (GCN-LSTM) was statistically important (p $<$ 0.05).

\begin{figure}[!t]
  \centering
  \includegraphics[width=\columnwidth]{Image3.jpg}
  \caption{Model performance evaluation. (Left) Mean Receiver Operating Characteristic (ROC) curve for the DST-GAT model across 5 folds, achieving a mean Area Under the Curve (AUROC) of 0.74. (Right) Training and validation loss curves for a representative fold, showing healthy training progression.}
  \label{fig:roc_loss_curves}
\end{figure}

\begin{table}[htbp]
\caption{Ablation studies for DST-GAT hyperparameters.}
\label{tab:ablation}
\begin{center}
\resizebox{\columnwidth}{!}{%
\begin{tabular}{@{}lc@{}}
\toprule
\textbf{Hyperparameter Variation} & \textbf{Mean AUC} \\ \midrule
\textbf{Sliding Window Size (W)} & \\
\quad W = 30 TRs (Stride=3) & 0.70 $\pm$ 0.04 \\
\quad \textbf{W = 50 TRs (Stride=5)} & \textbf{0.74 $\pm$ 0.03} \\
\quad W = 70 TRs (Stride=7) & 0.72 $\pm$ 0.03 \\ \midrule
\textbf{GAT Layers (Spatial Depth)} & \\
\quad L\_GAT = 1 & 0.68 $\pm$ 0.05 \\
\quad L\_GAT = 2 & 0.71 $\pm$ 0.04 \\
\quad \textbf{L\_GAT = 3} & \textbf{0.74 $\pm$ 0.03} \\
\quad L\_GAT = 4 & 0.73 $\pm$ 0.04 \\ \midrule
\textbf{Transformer Layers (Temporal Depth)} & \\
\quad L\_Trans = 2 & 0.69 $\pm$ 0.04 \\
\quad \textbf{L\_Trans = 4} & \textbf{0.74 $\pm$ 0.03} \\
\quad L\_Trans = 6 & 0.73 $\pm$ 0.04 \\ \bottomrule
\end{tabular}%
}
\end{center}
\small{Metrics are reported as mean $\pm$ standard deviation. Bold indicates the final selected value.}
\end{table}

The results are presented in Table \ref{tab:ablation}.

\subsection{Discovery of ASD Subtypes}
The identified ASD subtypes should be interpreted as exploratory, data-driven patterns derived from unsupervised clustering of learned representations. These subtype labels are not intended as clinically validated categories.

Using the final DST-GAT model trained on all data, we extracted embeddings for the 418 ASD subjects. Applying UMAP \parencite{mcinnes2018umap} for dimensionality reduction and HDBSCAN \parencite{campello2013density} for clustering to these embeddings revealed three distinct clusters, which we designate as Subtype 1, Subtype 2, and Subtype 3. The clustering result is visualized in Fig.~\ref{fig:umap_clusters}. HDBSCAN identified 354 subjects belonging to one of the three subtypes, while 64 subjects were classified as noise points (i.e., not clearly belonging to any of the cluster), potentially representing individuals with unique neurobiological profiles or a fourth, less-defined group. Post-hoc analysis of subtype-specific connectivity revealed significant differences in thalamo-cortical pathways \parencite{marko2015thalamo}, consistent with prior literature on heterogeneous neural substrates in ASD.

\begin{figure}[!t]
  \centering
  \includegraphics[width=\columnwidth]{Image4.jpg}
  \caption{UMAP visualization of the learned embeddings for the ASD cohort. HDBSCAN identified three distinct clusters, suggesting the presence of data-driven neurobiological subtypes. Subtype 1 (red, n=121), Subtype 2 (green, n=145), Subtype 3 (blue, n=88). Noise points are shown in grey (n=64).}
  \label{fig:umap_clusters}
\end{figure}


\begin{figure}[!t]
  \centering
  \includegraphics[width=\columnwidth]{Image5.jpg}
  \caption{Differential connectome maps for each identified ASD subtype compared to the TD group. Significant differences in functional connectivity patterns are highlighted and visualized spatially.}
  \label{fig:connectome_maps}
\end{figure}

% --- SECTION 5: DISCUSSION ---
% --- SECTION 5: DISCUSSION ---
\section{Discussion}
Our results demonstrate that modeling dynamic spatiotemporal brain connectivity leads to consistent performance improvements across evaluation metrics. However, beyond predictive accuracy, the primary value of the DST-GAT framework lies in its utility as a scientific tool for discovering data-driven neurobiological subtypes within the highly heterogeneous ASD population.

\subsection{Implications}
Clinical diagnostic criteria for ASD have historically been based on the behavioral presentations of specific demographic groups, contributing to documented disparities in diagnosis rates across sex and ethnicity. A recent meta-analysis found that the widely reported 4:1 male-to-female ratio is likely inflated by diagnostic bias, with population-based screening studies suggesting a ratio closer to 3:1 \parencite{loomes2017male}. Females and individuals from underrepresented populations who present with atypical profiles are at increased risk of delayed or missed diagnoses \parencite{lai2019sex}.

The unsupervised subtyping approach in our DST-GAT framework is relevant to this problem. Rather than relying solely on existing diagnostic labels, the model learns representations directly from neuroimaging data and uses clustering to identify subgroups that may not align with conventional clinical categories. The three neurobiological subtypes identified in this study, each characterized by distinct connectivity patterns in the default mode, salience, and fronto-parietal networks, suggest that ASD encompasses multiple separable neural profiles rather than a single homogeneous condition.

From a clinical standpoint, such data-driven stratification could inform more targeted intervention strategies. Current ASD therapies often follow a generalized approach, which may not account for the biological variability across individuals. If the subtypes identified here are validated in future independent cohorts, they could serve as a basis for matching patients to interventions that are better suited to their specific neurobiological profile, supporting a shift toward more personalized care in ASD.

\subsection{Limitations}
Importantly, the subtype structure identified in this study must be interpreted as an exploratory, data-driven pattern derived from learned representations. It has not yet been validated on independent, multi-site cohorts or correlated with longitudinal clinical and behavioral outcomes, which currently limits direct clinical deployment.

% --- SECTION 6: CONCLUSION ---
\section{Conclusion}
This study demonstrates that explicitly modeling spatiotemporal brain connectivity via the proposed DST-GAT framework improves ASD classification while enabling data-driven subgroup discovery. The identification of three distinct neurobiological subtypes within the ASD cohort, driven by differences in default mode, salience, and fronto-parietal connectivity, provides initial evidence that the heterogeneity of autism can be decomposed into separable neural profiles. Future work should focus on validating these subtypes in independent datasets and correlating them with clinical and behavioral outcomes to assess their translational potential.

\nocite{*}
\printbibliography

\end{document}
